\chapter{Ergodic theory}

Sometimes, we are gifted with the possibility to have neat formalization of dynamical systems, like in the case of the rotations $O(x)=\{x+nq, \mod 1\}$; topological dynamics was our first approach to look at orbits, looking at their topological properties, writing them in nice ways. 

Ergodic theory, instead, looks at dynamical systems form the probabilistic point of view: they are still deterministic systems, but the initial condition is probabilistic. To be more precise, ergodic theory deals with dynamical systems defined over a measure space $M$.Vaguely, the idea is that $\mu$ on M is a function $\mu: \mathcal{A\rightarrow \mathbb{R}^+_0}$, where $\mathcal{A}$ is the class of all measurable sets of M, and $\mu(A)$ is the measure of A. $\mu$ has a lot of natural properties, all stemming from \begin{enumerate}
    \item $\mu(\emptyset)=0$;
    \item if $\{A_n\}$ is a countable family of disjoint sets, then $\mu(\cup_n A_n)=\sum_{n}\mu(A_n)$;
\end{enumerate}
Every intuitive property (such as $\mu(A\cup B)=\mu(A)+\mu(B)-\mu(A\cap B)$) come from those. In practice, \begin{equation*}
\mu(A)=\int_A \phi(x) dLeb(x)|^{\text{with }\phi: M\rightarrow \mathbb{R}^+_0}
\end{equation*}
where dLeb(x) means using the Lebesgue measure.
Another possibility is to have the measure be \begin{equation*}
\mu=\delta_{\bar{x}}, \begin{cases}
    1, \bar{x}\in A \\ 0, \bar{x}\notin A
\end{cases}
\end{equation*}
so the delta function. This is used to measure a set, for example if you take a set A that does not contain $\bar{x}$, its measure is 0, while some set B that contains $\bar{x}$ has a measure of 1. You may easily see how this leads to the measure of a set by integrating it.